\subsubsection{静态调度结果与模型有效性校验}

完成联合优化模型和求解策略设计后，P2 还需要通过基础算例检验模型是否真正生成可执行的静态调度计划。这里的“可执行”指任务级计划同时满足任务分配、任务链衔接、候选路径选择、时间窗和电量安全等要求。

因此，本节不再重复讨论四类方法的横向优劣，而是围绕当前主流程的 m2 方案形成一条结果验收链：先说明代表性方案选择依据，再展示任务链和空间分布，最后从路径、电量、时间窗和负载均衡等方面检查模型约束是否落实到结果中。

\paragraph{代表性方案选择。}
前文给出了 m0、m1、m2、m3 四类求解策略。四类方法都可以生成静态任务链，但本节的重点不是进行方法横向比较，而是检验第 5.3 节联合优化模型和第 5.4 节分层目标规划能否形成可执行的基础计划。m2 是该模型体系的直接实现，能够在同一框架内同时处理任务分配、任务排序、候选路径选择、时间衔接和电量安全。因此，本节选取 m2 结果作为代表性算例；m0、m1 和 m3 分别作为基线对照、分解求解和规模扩展方法，其解质量与求解时间差异将在本章后续综合实验中统一比较。

\paragraph{静态任务链结果。}
表 \ref{tab:p2-static-sequence} 给出了 m2 在基础算例下得到的静态任务链。每一行对应一台 AMR，包括其执行任务序列、完成任务链后的时间以及最终剩余电量。

\begin{table}[!htbp]
\centering
\small
\begin{tabular}{p{0.14\linewidth}p{0.36\linewidth}p{0.18\linewidth}p{0.12\linewidth}}
\toprule
AMR & 任务序列 & 完成时间 & 剩余电量 \\
\midrule
AMR1 & T06 $\rightarrow$ T03 $\rightarrow$ T07 & 40.293444 & 54.7776 \\
AMR2 & T02 $\rightarrow$ T05 $\rightarrow$ T10 & 40.840111 & 48.0854 \\
AMR3 & T01 $\rightarrow$ T09 $\rightarrow$ T12 & 38.944691 & 67.5520 \\
AMR4 & T04 $\rightarrow$ T08 $\rightarrow$ T11 & 38.596970 & 43.4962 \\
\bottomrule
\end{tabular}
\caption{P2 基础算例静态任务链结果}
\label{tab:p2-static-sequence}
\end{table}

可以看到，12 个任务被均匀分配给 4 台 AMR，每台 AMR 执行 3 个任务。四台 AMR 的完成时间集中在 38.60 到 40.84 之间，系统完工时间为 40.840111，说明模型没有将任务集中压给少数车辆，而是在完工时间和负载均衡之间形成了较平衡的计划。最低剩余电量为 43.4962，也说明电量安全约束参与了任务链选择，而不是在结果生成后再进行事后检查。

\paragraph{空间分布校验。}
仅从任务链表可以看出任务归属和完成时间，但还不能直观看出任务链在仓库平面中的分布。图 \ref{fig:p2-static-map} 将表 \ref{tab:p2-static-sequence} 的结果投影到二维仓库场景中，其中虚线表示任务之间的空驶衔接，实线表示载货运输，圆形序号表示同一 AMR 内部的执行顺序。

\begin{figure}[!htbp]
\centering
\includegraphics[width=\linewidth]{p2_assignment_map_m2.png}
\caption{P2 基础算例任务链空间分布}
\label{fig:p2-static-map}
\end{figure}

图 \ref{fig:p2-static-map} 表明，P2 结果不是简单的任务平均分摊，而是在空间位置、任务时间窗和 AMR 起点共同作用下形成四条相对紧凑的任务链。例如，每台 AMR 的任务顺序均从其初始位置出发，并通过空驶衔接连接到后续取货点；载货段则连接取货点与对应送货点。这说明第 5.3 节中的任务链约束和时间衔接约束已经体现在空间执行顺序中。

\paragraph{可行性与接口校验。}
为进一步检验模型约束是否真正起作用，本文从路径可达性、电量安全、时间窗服务质量、系统完工时间和负载均衡五个方面对结果进行校验，如表 \ref{tab:p2-feasibility-check} 所示。

\begin{table}[!htbp]
\centering
\small
\begin{tabular}{p{0.22\linewidth}p{0.20\linewidth}p{0.42\linewidth}}
\toprule
校验内容 & 结果 & 说明 \\
\midrule
空驶路径 & 0 处缺失 & 每个任务前的衔接移动都有可用候选路径 \\
载货路径 & 0 处缺失 & 每个取货到送货过程都有可用候选路径 \\
电量安全 & 0 次违反 & 任务链执行后未低于安全电量阈值 \\
时间窗 & 0 个延期任务 & 静态计划满足任务时间窗要求 \\
系统完工时间 & 40.840111 & 四条任务链均在该时刻前完成 \\
负载均衡 & 约 2.24 & 四台 AMR 完成时间差较小 \\
\bottomrule
\end{tabular}
\caption{P2 基础算例可行性验收}
\label{tab:p2-feasibility-check}
\end{table}

表 \ref{tab:p2-feasibility-check} 说明，第 5.3 节提出的路径可达性、任务链时间衔接、时间窗和电量安全约束在基础算例中均得到满足；完成时间差约为 2.24，也呼应了第 5.4 节中将 $C_{\max}$ 和负载均衡纳入目标规划的设计。由此可见，P2 输出的不只是任务归属，而是一份包含任务顺序、路径选择、时间估算和电量状态的静态任务级计划。
